% ID: SEL-D0023-OA-E04
% ESTADO: cotejada. Segunda pasada visual completa y borrador compilado revisados.
% Fuente: Selectividad/matematicas_aplicadas/reserva_b_EXAMEN 2 -ANDALUCÍA.pdf
% SHA-256: 8cdb86259405652622f2ea129bffbdb239076ac612fd7fcebb86b9d5bdc451ec
% Página: 1; opción A; ejercicio original 4.

Se sabe que el diámetro de las estrellas de mar de una región sigue una ley Normal con varianza $2.25\ \text{cm}^2$. Se sospecha que, igual que ocurre en otras regiones, su diámetro no supera los $11.7\ \text{cm}$ $(H_0:\mu\leq 11.7)$. Para confirmarlo se extrae una muestra aleatoria de estrellas de mar de esa región, obteniéndose los siguientes diámetros:
\[
12.5\quad11.8\quad13.1\quad14.3\quad11.7\quad12.6\quad12.7\quad12.1\quad13.5\quad11.5
\]
\begin{enumerate}[label=\alph*)]
\item \textbf{(1.75 puntos)} Plantee un contraste de hipótesis, y para un nivel de significación del 5\%, obtenga la región de rechazo del contraste. \textquestiondown{}Se puede confirmar la sospecha?
\item \textbf{(0.75 puntos)} \textquestiondown{}Y para un nivel de significación del 3\%, se puede confirmar la sospecha?
\end{enumerate}
